\documentclass[11pt,a4paper]{article}
\usepackage[T1]{fontenc}
\usepackage{isabelle,isabellesym}
\usepackage{pdfsetup}
\urlstyle{rm}
\isabellestyle{it}

\begin{document}

\title{The Schur Number $S(4) = 44$, via Kernel Replay of LRAT Proofs}
\author{Carles Mar\'in}
\maketitle

\begin{abstract}
We prove in Isabelle/HOL that the Schur number $S(4)$ is 44: $\{1,\dots,44\}$ has a 4-colouring without
monochromatic solutions of $a+b=c$, and $\{1,\dots,45\}$ has none. The statement is about colourings. A colour
symmetry, proved in HOL, reduces the upper bound to four cases; each is encoded by explicit instances of three
facts about the colouring, refuted by a SAT solver, and its LRAT proof is replayed by the Isabelle kernel. The
replay method is part of the entry: it accepts RUP steps only, turns each into a resolution chain and replays
them forwards, dropping every clause theorem after its last use. No oracle is used.
\end{abstract}

\tableofcontents

\parindent 0pt\parskip 0.5ex

\input{session}

\end{document}
